% Consolidated guide-selftest03 from 97b62cc61307eba158a4498b55de3a5fdef3ff43; source fragments preserved.

% === BEGIN guide-selftest03-choices.tex blob d9b683c6946f0ed0217921c1a2ef2be419ccc866 ===
% 第15批：《大学物理学学习指导》第5版自测题3“相对论”
% 选择3.1.1—3.1.10，原书印刷54—56页／PDF物理62—64页。
\subsection{自测题3：速度合成、时钟与相对论能量选择题}

\sourcetest{3.1.1}{54}{62}
\problemtext 长度为$L$的火筒匀速行驶，
人从其后端沿筒射出子弹，子弹\emph{相对火筒}的速率为$v_2$，
问火筒自有时钟测得子弹到前靶的时间。
\workings
筒身自身静止系中两端相距固有长度$L$，
子弹相对筒速率已知$v_2$，
\[
\boxed{\Delta t_{\rm tube}=\frac L{v_2}}.
\]
选\textbf{B}。
\crosscheck 不应将火筒对地速率$v_1$与子弹对筒速率
简单相加用于\emph{筒自身时钟}的传播时间。

\sourcetest{3.1.2}{54}{62}
\problemtext 飞船在地面系以$v$飞行，船首向船尾发送光信号，
飞船上接收者的本地时钟记录两事件相隔$\Delta t$，
求飞船固有长度。
\workings
船自身系两端静止，真空光速$c$，故
\[
\boxed{L_0=c\Delta t}.
\]
选\textbf{A}。
\crosscheck 信号发射在船首、接收在船尾，
两事件不是同地，故$\Delta t$不是这对事件的严格固有时，
虽然船的\emph{长度}是固有长度。

\sourcetest{3.1.3}{54}{62}
\problemtext 直尺静止于$K'$系，与$x'$轴成45°；
$K'$以非零速率沿$x$正向相对$K$运动，
求$K$中尺的夹角相对45°是增大还是减小。
\workings
同步测长时，沿运动方向长度缩短$\Delta x=\Delta x'/\gamma$，
垂直方向不变$\Delta y=\Delta y'$，
\[
\tan\theta=\frac{\Delta y}{\Delta x}
=\gamma\tan45^\circ=\gamma>1.
\]
故$\boxed{\theta>45^\circ}$，选\textbf{A}。
\crosscheck 夹角改变与运动方向的\emph{正负号}无关，只与$\gamma$有关。

\sourcetest{3.1.4}{54--55}{62--63}
\problemtext 判断：一惯性系中同地同时发生的两事件，
及另一对同一时刻却不同地的两事件，
在另一运动惯性系中是否仍可同时。
\workings
同地同时就是同一时空事件，
其$\Delta t=\Delta x=0$在任意惯性变换后仍为0。
对\emph{不同地}但同一时刻的事件，
沿分离方向作一维洛伦兹变换
$\Delta t'=\gamma(0-v\Delta x/c^2)\ne0$
（$v\Delta x\ne0$），
故通常失去同时性。按原题默认的一维相对运动约定，
选\textbf{A（第一组同时、第二组不同）}。
\crosscheck 若两个事件的三维空间位移
严格垂直于参考系相对速度，
仍有$\bm v\cdot\Delta\bm r=0$而继续同时。
所以该选择结论需要\emph{沿事件分离方向有非零速度分量}的约定。

\sourcetest{3.1.5}{55}{63}
\problemtext 电子静能约$0.5$MeV，动能约$0.25$MeV，
求其速率与光速的比。
\workings
$(\gamma-1)m_ec^2=0.25$MeV，$m_ec^2=0.50$MeV，
故$\gamma=1.5$，
\[
\boxed{\beta=\frac vc=\sqrt{1-\frac1{1.5^2}}
=\frac{\sqrt5}{3}\simeq0.745}.
\]
选\textbf{C（约$0.75c$）}。
\crosscheck 经典式$K=\tfrac12mv^2$在如此高能下不可使用。

\sourcetest{3.1.6}{55}{63}
\problemtext 关于狭义相对论的四条判断：
有静质量粒子速度限制、测量时间长度依参考系变化、
同时性的绝对性、相对运动时钟的时间延缓。
\workings
有静质量物体不能加速达到或超过$c$，命题(1)真；
长度和时间的比较可能依惯性系改变，(2)真；
分离两事件的同时性不是绝对的，(3)假；
相对运动钟的坐标时间间隔相对固有时可延长，(4)真。
选\textbf{B（1、2、4）}。
\crosscheck 原书术语“质量随运动状态改变”应在现代写法中
区分静质量不变与总能量$E=\gamma m_0c^2$可变。

\sourcetest{3.1.7}{55}{63}
\problemtext 飞船自身长度90m，相对地球以$0.8c$运动。
船尾发射光脉冲并于船头接收到，
求地面系中这两个事件的空间距离。
\workings
地面测得飞船长度
$L=L_0\sqrt{1-0.8^2}=90(0.6)=54$m。
在地面系，光追赶移动船头：
$c\,t=L+0.8c\,t$，
故$ct=54/(1-0.8)=270$m。
\[
\boxed{\Delta x_{\rm Earth}=270\,\mathrm m}.
\]
选\textbf{C}。
\crosscheck 题问的是\emph{发射事件到接收事件的空间间隔}，
并非地面同一时刻测量的飞船长度54m。

\sourcetest{3.1.8}{55}{63}
\problemtext 电子以速率$0.99c$运动，静能取0.51MeV，
求其相对论动能。
\workings
\[
\boxed{K=(\gamma-1)m_ec^2
=\left(\frac1{\sqrt{1-0.99^2}}-1\right)(0.51)
\simeq3.10\,\mathrm{MeV}}.
\]
选\textbf{C}。
\crosscheck 粒子速度接近光速时，经典动能公式严重低估。

\sourcetest{3.1.9}{55}{63}
\problemtext 年发电量为100亿千瓦时的核电站，
若电能完全来自核材料静质量减少，
估算所对应的质量亏损。
\workings
100亿kWh=$10^{10}$kWh，换算能量
$E=3.6\times10^{16}$J。
按$E=\Delta mc^2$，
\[
\boxed{\Delta m=\frac{3.6\times10^{16}}
{(3.0\times10^8)^2}=0.40\,\mathrm{kg}}.
\]
选\textbf{A}。
\crosscheck 这是假设100%质能转化为电能的理想质量差，
并非实际消耗的燃料总质量。

\sourcetest{3.1.10}{55--56}{63--64}
\problemtext 两个静质量均为$m_0$的粒子在实验室系
沿一条直线以速度$\pm v$相向运动，
完全非弹性碰合为一个静止复合粒子，
求它的静质量$M_0$。
\workings
初始总动量零，初总能量为$2\gamma m_0c^2$；
碰合后的复合物静止，故
$M_0c^2=2\gamma m_0c^2$，
\[
\boxed{M_0=\frac{2m_0}{\sqrt{1-v^2/c^2}}}.
\]
选\textbf{D}。
\crosscheck 与\emph{同速同向}合并仅得$2m_0$不同；
这里初始的相对动能转化成复合粒子的内能，
增加系统静质量。
% === END guide-selftest03-choices.tex ===

% === BEGIN guide-selftest03-blanks.tex blob db255392882f30fc41ecbafe7b2bf0199698fd39 ===
% 第15批：《大学物理学学习指导》第5版自测题3填空3.2.1—3.2.8
% 印刷56／PDF物理64页。
\subsection{自测题3：光速、时间延缓和相对论密度填空}

\sourcetest{3.2.1}{56}{64}
\problemtext 恒星相对地球以速率$v$匀速运动；
恒星辐射的光子相对地球的速率大小是多少？
\workings
狭义相对论的真空光速不变，
\[
\boxed{u_\gamma=c}.
\]
\crosscheck 不能用伽利略相加写成$c+v$或$c-v$。

\sourcetest{3.2.2}{56}{64}
\problemtext $S'$沿$S$系$x$轴\emph{负方向}
以$0.5c$运动；$S'$原点O'沿$x'$正向发出光波。
$S$测量该光速大小。
\workings
相对运动速度$v=-0.5c$，
速度合成
\[
u=\frac{c-0.5c}{1-(0.5)}=c,
\]
故$\boxed{u=c}$。
\crosscheck 对光速使用洛伦兹速度合成；
不能得到经典的$0.5c$。

\sourcetest{3.2.3}{56}{64}
\problemtext 不稳定$\pi^+$介子的自身平均寿命为
$2.6\times10^{-8}$s，以$0.8c$相对实验室运动，
求实验室平均寿命。
\workings
运动钟时间延缓
\[
\boxed{\tau_{\rm lab}
=\frac{\tau_0}{\sqrt{1-0.8^2}}
=\frac{2.6\times10^{-8}}{0.6}
\simeq4.33\times10^{-8}\,\mathrm s}.
\]
\crosscheck 这是\emph{平均寿命}的参考系变换，
不是每个粒子实际存活时间都等于该值。

\sourcetest{3.2.4}{56}{64}
\problemtext 两位观察者相向接近的相对速度为$0.6c$，
其中O测得初始距离20m。
求O'自己的时钟从该初始时刻到相遇所历时间。
\workings
在O静止系，两者相遇需
$t_O=20/(0.6c)$。
O'时钟测量自身世界线上两个事件的固有时间，
\[
\boxed{\Delta\tau_{O'}=
t_O\sqrt{1-0.6^2}
=\frac{20(0.8)}{0.6c}
\simeq8.89\times10^{-8}\,\mathrm s}.
\]
\crosscheck 20m距离是由O系同时量得的初始距离；
不能直接套O'系长度20m而不考虑收缩。

\sourcetest{3.2.5}{56}{64}
\problemtext (1) 粒子的相对论动量大小何时等于相同速率下的
经典动量大小的2倍？
(2) 相对论动能何时等于自身静能？
\workings
(1) $p_{\rm rel}=\gamma m_0v$，
$p_{\rm classical}=m_0v$，要求$\gamma=2$。
(2) $K=(\gamma-1)m_0c^2=m_0c^2$，
也要求$\gamma=2$，所以两问速度相同：
\[
\boxed{v=c\sqrt{1-\tfrac1{2^2}}
=\frac{\sqrt3}{2}c}.
\]
\crosscheck 这里按静质量$m_0$定义经典动量，
不是沿用运动质量$m_{\rm rel}=\gamma m_0$后又多乘一个$\gamma$。

\sourcetest{3.2.6}{56}{64}
\problemtext 电子由静止加速到$0.6c$，求外力所做功，
以电子静质量$m_0$表达。
\workings
$\gamma=(1-0.6^2)^{-1/2}=1.25$，
\[
\boxed{W=\Delta K=(\gamma-1)m_0c^2
=\tfrac14m_0c^2}.
\]
\crosscheck 经典值$\tfrac12 m_0(0.6c)^2=0.18m_0c^2$，
与所求相对论值$0.25m_0c^2$不同。

\sourcetest{3.2.7}{56}{64}
\problemtext 观察者甲相对乙以$4c/5$运动，
甲携带静止长度$l$、截面积$S$、静止质量$m$的细棒，
细棒长度沿运动方向。
求甲、乙各自记录的棒的密度。
\workings
甲在棒自身静止系中测得
\[
\boxed{\rho_{\rm rest}=\frac{m}{Sl}}.
\]
乙测量的同时长度为
$l_B=l\sqrt{1-(4/5)^2}=3l/5$。
若使用\emph{现代静质量密度}的定义，
质量仍为$m$，故
\[
\boxed{\rho_{B,\rm restmass}=\frac{5m}{3Sl}}.
\]
但源书时代部分章节将“运动质量”
定义为$m_{\rm rel}=\gamma m=(5/3)m$；
若题中“密度”意为\emph{运动质量除同时体积}，
则乙得到
\[
\boxed{\rho_{B,\rm relativisticmass}
=\frac{(5/3)m}{(3/5)Sl}=\frac{25m}{9Sl}}.
\]
\crosscheck \textbf{源题术语约定必须独立二校：}
两种密度的数值不同，并非同一物理定义的两个解。
现代连续介质密度和相对论能动张量的定义又需考虑参考系切片；
本题只登记两种明确的教材简化约定，不虚称唯一无条件值。

\sourcetest{3.2.8}{56}{64}
\problemtext 速率为$0.99c$的电子静质量为
$9.11\times10^{-31}$kg。求相对论总能量（J）
及经典动能与相对论动能之比。
\workings
设$\gamma=1/\sqrt{1-0.99^2}\simeq7.0888$，
则
\[
\boxed{E=\gamma m_0c^2
\simeq(7.0888)(9.11\times10^{-31})(3.0\times10^8)^2
\simeq5.81\times10^{-13}\,\mathrm J}.
\]
经典动能与相对论动能分别是
$K_{\rm cl}=\tfrac12m_0v^2$、
$K_{\rm rel}=(\gamma-1)m_0c^2$，所以
\[
\boxed{\frac{K_{\rm cl}}{K_{\rm rel}}
=\frac{0.5(0.99)^2}{\gamma-1}
\simeq0.0805}.
\]
\crosscheck “总能量”包括静能，不是相对论动能本身。
% === END guide-selftest03-blanks.tex ===

% === BEGIN guide-selftest03-solutions.tex blob a132cfd051b1177f1ae694ff9904ca4e28587a03 ===
% 自测题3计算3.3.1—3.3.7与问答3.4.1；
% 原书印刷56—57页／PDF物理64—65页，特别注明同时性与能量定义。
\subsection{自测题3：洛伦兹变换、时间间隔和质能}

\sourcetest{3.3.1}{56}{64}
\problemtext 甲、乙静止于惯性系$K,K'$，
甲测得两个同地事件的时间间隔为4s，
乙测得对应间隔为5s。
求$K'$相对$K$的速率，及乙测得这两个事件的空间距离。
\workings
甲系两事件同地，故4s为固有时间。
$\gamma=5/4$，于是
\[
\boxed{v=c\sqrt{1-(4/5)^2}=0.60c}.
\]
若$K'$沿$+x$相对$K$移动，
$\Delta x=0$且$\Delta t=4$s，
由洛伦兹变换
\[
\boxed{\Delta x'=-\gamma v\Delta t
=-(5/4)(0.6c)(4\,\mathrm s)
=-3c\,\mathrm s\simeq-9.0\times10^8\,\mathrm m}.
\]
空间\emph{间隔大小}为$9.0\times10^8$m。
\crosscheck 原题只给相对速率大小，
因此有向位移正负依选定坐标轴；不能仅依据两时间读出唯一方向。

\sourcetest{3.3.2}{57}{65}
\problemtext 飞船固有长度$L_0=90$m，
对地速度$0.8c$，掠过地面一固定观测站。
求(1)观测站地面时钟记录整艘飞船通过的时间，
(2)飞船本地系测得观测站从船首到船尾经过的时间。
\workings
地面同时测长$L=L_0\sqrt{1-0.8^2}=54$m，
所以地面观测站上同地两事件所隔的固有时间
\[
\boxed{\Delta t_{\rm ground}=
\frac L{0.8c}=\frac{54}{0.8c}
\simeq2.25\times10^{-7}\,\mathrm s}.
\]
飞船系中本船静止长度为90m，
地面站以$0.8c$穿过两端，故
\[
\boxed{\Delta t_{\rm ship}=
\frac{L_0}{0.8c}
\simeq3.75\times10^{-7}\,\mathrm s}.
\]
\crosscheck 船首遇站和船尾遇站在\emph{地面同地}、
在\emph{飞船异地}，因此两系时间间隔不相同；
$\Delta t_{\rm ship}=\gamma\Delta t_{\rm ground}$。

\sourcetest{3.3.3}{57}{65}
\problemtext 在$K$系中两个事件沿$x$的间距
$\Delta x=500$m、时间差$\Delta t=2\times10^{-7}$s。
若$K'$系观察者测得这两个事件同时，
求$K'$沿$x$相对$K$的有向速度。
\workings
$\Delta t'=\gamma(\Delta t-v\Delta x/c^2)=0$，
从而
\[
\boxed{v=\frac{c^2\Delta t}{\Delta x}
=\frac{(3.0\times10^8)^2(2\times10^{-7})}{500}
=3.6\times10^7\,\mathrm{m/s}=0.12c}.
\]
\crosscheck $|\Delta x|>c|\Delta t|$，事件为空间型间隔，
存在将其视为同时的惯性系；若是类时分离则不可如此变换。

\sourcetest{3.3.4}{57}{65}
\problemtext 地面发射台左右的西站W、东站E距台均为$L_0$。
发射台同一时刻向两站发出同一光信号。
一架飞机以$v$沿从W到E的方向匀速飞过，
求飞机上记录两接收站\emph{收到同一信号}
这两个事件的时间间隔。
\workings
发射台与两站均静止于地面系$K$。
两接收事件在地面系同时：
$t_E=t_W=L_0/c$，位置
$x_E=+L_0$、$x_W=-L_0$，
所以$\Delta t=0$、$\Delta x=2L_0$。
飞机系沿$+x$相对地面以$v$运动，
\[
\Delta t'=t_E'-t_W'
=\gamma\left(\Delta t-\frac{v\Delta x}{c^2}\right)
=-\frac{2\gamma vL_0}{c^2}.
\]
故
\[
\boxed{|\Delta t'|=
\frac{2vL_0}{c^2\sqrt{1-v^2/c^2}}}.
\]
负号表明飞机系认为\textbf{东站先接收}。
\crosscheck 切不可把两接收站误当成\emph{随飞机运动}，
也不可用$L_0/(c-v)$描述地面静止站的接收时刻。

\sourcetest{3.3.5}{57}{65}
\problemtext 一宇宙射线介子动能
$K=7M_0c^2$，
其中$M_0$是其静止质量，
求实验室系测得寿命与介子固有平均寿命之比。
\workings
$K=(\gamma-1)M_0c^2=7M_0c^2$，
故$\gamma=8$。
时间延缓使
\[
\boxed{\tau_{\rm lab}/\tau_0=8}.
\]
\crosscheck 不可直接把动能是7倍静能误写为时间延缓因子7；
延缓因子是总能量与静能之比8。

\sourcetest{3.3.6}{57}{65}
\problemtext 将电子速率从$1.2\times10^8$m/s
提高至$2.4\times10^8$m/s，
求所需外力的相对论净功，
取电子静质量$9.11\times10^{-31}$kg。
\workings
采用$c\simeq3.0\times10^8$m/s，
$\beta_1=0.4,\beta_2=0.8$，
$\gamma_1=1/\sqrt{0.84}$、
$\gamma_2=1/\sqrt{0.36}=5/3$。
相对论功—能关系
\[
\begin{aligned}
W&=(\gamma_2-\gamma_1)m_0c^2\\
&=\left(\frac53-\frac1{\sqrt{0.84}}\right)
(9.11\times10^{-31})(3.0\times10^8)^2,
\end{aligned}
\]
故
\[
\boxed{W\simeq4.72\times10^{-14}\,\mathrm J
\simeq0.294\,\mathrm{MeV}}.
\]
\crosscheck 只改变动能而不改变静能，
故总能量差与动能差相同。

\sourcetest{3.3.7}{57}{65}
\problemtext 观察者甲以$0.8c$相对静止的乙运动，
甲携带静止质量1kg的物体。
两人各自测得该物体的\emph{总能量}是多少？
\workings
物体在甲系静止，所以
\[
\boxed{E_{\rm A}=mc^2
\simeq8.99\times10^{16}\,\mathrm J}.
\]
乙系中物体以$0.8c$运动，
$\gamma=1/\sqrt{1-0.8^2}=5/3$，
\[
\boxed{E_{\rm B}=\gamma mc^2
=\tfrac53mc^2
\simeq1.50\times10^{17}\,\mathrm J}.
\]
\crosscheck 总能量不是洛伦兹标量，
\emph{静质量}才是；甲不能看到自己携带的物体有额外平动动能。

\subsection{自测题3：经典与相对论不变量问答}

\sourcetest{3.4.1}{57}{65}
\problemtext 对位移、质量、时间、速度、动能这五个物理量，
比较其在经典伽利略变换与狭义相对论洛伦兹变换下的表达，
指出哪些可保持不变。
\workings
一维伽利略变换为
$x'=x-vt,\ t'=t$，速度$u'=u-v$；
静质量$m'=m$、经典动能
$K'=\tfrac12m(u-v)^2$，
故\textbf{时间差和质量}一般保持不变；
\emph{同一时刻两点的空间距离}也保持不变，
但不同时间发生的两事件位移差
$\Delta x'=\Delta x-v\Delta t$并非一般不变量。
洛伦兹变换为
\[
x'=\gamma(x-vt),\qquad
t'=\gamma(t-vx/c^2).
\]
长度/位移、坐标时间、速度、动能一般随惯性系改变，
现代定义的\textbf{静质量$m_0$保持不变}，
同时有时空间隔
\[
\boxed{c^2\Delta t'^2-\Delta x'^2
=c^2\Delta t^2-\Delta x^2}
\]
和单粒子质能关系
$\boxed{E^2-(pc)^2=m_0^2c^4}$保持不变。
\crosscheck 若“质量”沿用原书旧式\emph{运动质量}
$m_{\rm rel}=\gamma m_0$，则它随速度变，不是洛伦兹标量；
因此应事先说明所说的是现代\emph{静质量}
还是教材历史术语的\emph{运动质量}。
% === END guide-selftest03-solutions.tex ===
